\begin{abstract}
This thesis investigates the predictive value of sentiment extracted from financial filings for stock market forecasting across multiple modeling settings, aggregation levels, and temporal architectures. 
Using FinBERT-derived sentiment features, we evaluate classification and regression tasks under company, sector, and industry-level grouping, as well as temporal modeling with LSTM and transformer-based PatchTST architectures. 
The results show that sentiment information is consistently more informative for directional prediction than for numerical return estimation. 
For classification, we observe statistically significant improvements of up to 7.39\% in average F1 score, particularly at the sector aggregation level and when temporal modeling is employed. 
In contrast, regression performance remains more challenging, although transformer-based temporal models achieve statistically significant reductions in RMSE of up to -11.35\% on average. 
At the individual entity level, sentiment signals can produce substantially larger effects, with improvements of up to 169.29\% in F1 score for classification and decreases of up to -53.13\% in RMSE for regression. 
These findings highlight the importance of aggregation level, label construction, and temporal modeling in extracting predictive value from financial text, and suggest that sentiment derived from corporate disclosures is most effective for directional forecasting while numerical return prediction remains more model-dependent.
\end{abstract}